\scanpage{15}
\mspage{13}
e) Le complexe de chaînes $C_\bullet(\unsure{\mathscr X|X|\mathcal X}_\ast,\deleted{\unsure{\mathbb Z|k}}k)$ permet de réconstruire tous les complexes de cochaînes $C^\bullet(\mathscr X_\ast,k')$, pour $k'$ \deleted{variable} \inserted{$k$-alg.}, par
\[
C^\bullet(\mathscr X_\ast,k)\simeq\mathrm{Hom}_{\mathbb Z}(C_\bullet(\mathscr X_\ast,\deleted{\unsure{\mathbb Z}}k),k).
\]
\inserted{(L'objet ``le plus fin'' est du $C_\bullet(\mathscr X_\ast,\mathbb Z)$, on a alors $C_\bullet(\mathscr X_\ast,k)\simeq C_\bullet(\mathscr X_\ast,\mathbb Z)\otimes_{\mathbb Z} k$)}

Il est possible de \deleted{\unsure{au}} de définir une cobigèbre \inserted{diff} $C_{\bullet\bullet}^{\mathrm{DRpd}}(\mathscr X_\ast,\deleted{\unsure{\mathbb Z|k}}k)$ \inserted{$k$-coaugmentée} à copuiss. div, telle que l'on ait, pour tte $k^{+}$-alg. $k'$
\[
C^{\bullet\bullet}_{\mathrm{DRpd}}(\mathscr X_\ast,k')\simeq\mathrm{Hom}_k(C_{\bullet\bullet}^{\mathrm{DRpd}}(\mathscr X_\ast,k),k')
\]
(compatible avec ttes les structures).

On aura d'ailleurs
\[
C_{\bullet\bullet}^{\mathrm{DRpd}}(\mathscr X_\ast,k')\simeq C_{\bullet\bullet}^{\mathrm{DRpd}}(\mathscr X_\ast,k)\otimes_k k'
\]
L'objet le plus fin est $C_{\bullet\bullet}^{\mathrm{DRpd}}(\mathscr X_\ast,\mathbb Z)$. C'est lui qu'il conviendrait de considérer (au lieu de son ``dual'' $C^{\bullet\bullet}_{\mathrm{DRpd}}(\mathscr X_\ast,\mathbb Z)$) si on veut aborder b) c) d) sans conditions de finitude.

Notons que [tt comme $C^{\bullet\unsure{\ast|\bullet}}_{\mathrm{DRpd}}(\mathscr X_\ast,k)$, pour $\mathscr X_\ast$ variable, transforme $\underset{\longrightarrow}{\lim}$ quelq. en $\underset{\longleftarrow}{\lim}$], $C_{\bullet\bullet}^{\mathrm{DRpd}}(\mathscr X_\ast,\mathbb Z)$ \deleted{transforme} \inserted{commute aux} $\underset{\longrightarrow}{\lim}$ quelconques (pas néc. filtrantes)

% archive sheet number: 14
